\section{AHSME 1969}
        \begin{enumerate}
        	\item (\textit{AHSME 1969}) When $x$ is added to both the numerator and denominator of the fraction 
$\frac{a}{b},a \ne b,b \ne 0$, the value of the fraction is changed to $\frac{c}{d}$. 
Then $x$ equals:


 $\text{(A) } \frac{1}{c-d}\quad \text{(B) } \frac{ad-bc}{c-d}\quad \text{(C) } \frac{ad-bc}{c+d}\quad \text{(D) }\frac{bc-ad}{c-d} \quad \text{(E) } \frac{bc+ad}{c-d}$

 \textbf{Answer: }B



\textbf{Solution}

$\frac{a+x}{b+x}=\frac{c}{d}$,

 $bc+cx=ad+dx$,

 $(c-d)x=ad-bc$,

 $x=\frac{ad-bc}{c-d}$. The answer is $\fbox{B}$.

 


	\item (\textit{AHSME 1969}) If an item is sold for $x$ dollars, there is a loss of $15\%$ based on the cost. If, however, the same item is sold for $y$ dollars, there is a profit of $15\%$ based on the cost. The ratio of $y:x$ is:


 $\text{(A) } 23:17\quad \text{(B) } 17y:23\quad \text{(C) } 23x:17\quad \\ \text{(D) dependent upon the cost} \quad \text{(E) none of these.}$

 \textbf{Answer: }A



\textbf{Solution}

Let $c$ be the cost. Selling the item for $x$ dollars equates to losing $15\%$ of $c$, so $x=.85c$. Selling the item for $y$ dollars equates to profiting by $15\%$ of $c$, so $y=1.15c$. Therefore $\frac{y}{x}=\frac{1.15c}{.85c}=\frac{23}{17}$. The answer is $\fbox{A}$.

 


	\item (\textit{AHSME 1969}) If $N$, written in base $2$, is $11000$, the integer immediately preceding $N$, written in base $2$, is:


 $\text{(A) } 10001\quad \text{(B) } 10010\quad \text{(C) } 10011\quad \text{(D) } 10110\quad \text{(E) } 10111$

 \textbf{Answer: }E



\textbf{Solution}

$11000_2-1_2=10111_2$. $\fbox{E}$



	\item (\textit{AHSME 1969}) Let a binary operation $\star$ on ordered pairs of integers be defined by $(a,b)\star (c,d)=(a-c,b+d)$. Then, if $(3,3)\star (0,0)$ and $(x,y)\star (3,2)$ represent identical pairs, $x$ equals:


 $\text{(A) } -3\quad \text{(B) } 0\quad \text{(C) } 2\quad \text{(D) } 3\quad \text{(E) } 6$

 \textbf{Answer: }E



\textbf{Solution}

Performing the operation based on the definition, $(3,3)\star(0,0) = (3,3)$ and $(x,y)\star(3,2)=(x-3,y+2)$.  Because the outputs are identical pairs, they must equal each other, so $3 = x-3$.  Solving for x yields $x = 6$, which is answer choice $\boxed{\textbf{(E)}}$.

 


	\item (\textit{AHSME 1969}) If a number $N,N \ne 0$, diminished by four times its reciprocal, equals a given real constant $R$, then, for this given $R$, the sum of all such possible values of $N$ is


 $\text{(A) } \frac{1}{R}\quad \text{(B) } R\quad \text{(C) } 4\quad \text{(D) } \frac{1}{4}\quad \text{(E) } -R$

 \textbf{Answer: }B



\textbf{Solution}

Write an equation from the given information.
 \[N - \frac{4}{N} = R\]
 \[N^2 - 4 = RN\]
 \[N^2 - RN - 4 = 0\]By \href{https://artofproblemsolving.com/wiki/index.php?title=Vieta\%27s\_Formulas}{Vieta's Formulas}, the sum of all possible values of $N$ for a given $R$ is $R$, so the answer is $\boxed{\textbf{(B)}}$.

 


	\item (\textit{AHSME 1969}) The area of the ring between two concentric circles is $12\frac{1}{2}\pi$ square inches. The length of a chord of the larger circle tangent to the smaller circle, in inches, is:


 $\text{(A) } \frac{5}{\sqrt{2}}\quad \text{(B) } 5\quad \text{(C) } 5\sqrt{2}\quad \text{(D) } 10\quad \text{(E) } 10\sqrt{2}$

 \textbf{Answer: }C



\textbf{Solution}


                    \begin{figure}[H]
                        \centering
                        \includegraphics[width=0.5\textwidth]{Solution/AHSME_I_1969_6_Solution_Figure_1.png}
                    \end{figure}
                    Let $a$ be radius of larger circle, and $b$ be radius of smaller circle.  The area of the ring can be determined by subtracting the area of the smaller circle from the area of the larger circle, so
 \[\pi a^2 - \pi b^2 = \frac{25 \pi}{2}\]
 \[a^2 - b^2 = \frac{25}{2}\]From the diagram, by using the \href{https://artofproblemsolving.com/wiki/index.php?title=Pythagorean\_Theorem}{Pythagorean Theorem}, half of the chord is $\frac{5}{\sqrt{2}} = \frac{5 \sqrt{2}}{2}$ units long, so the the length of the entire chord is $\boxed{\textbf{(C) } 5 \sqrt{2}}$.

 


 


	\item (\textit{AHSME 1969}) If the points $(1,y_1)$ and $(-1,y_2)$ lie on the graph of $y=ax^2+bx+c$, and $y_1-y_2=-6$, then $b$ equals:


 $\text{(A) } -3\quad \text{(B) } 0\quad \text{(C) } 3\quad \text{(D) } \sqrt{ac}\quad \text{(E) } \frac{a+c}{2}$

 \textbf{Answer: }A



\textbf{Solution}

Because the two points are on the quadratic, $y_1 = a + b + c$ and $y_2 = a - b + c$.  Because $y_1 - y_2 = -6$,
 \[(a+b+c)-(a-b+c)=-6\]
 \[2b=-6\]
 \[b=-3\]The answer is $\boxed{\textbf{(A)}}$.

 


	\item (\textit{AHSME 1969}) Triangle $ABC$ is inscribed in a circle. The measure of the non-overlapping minor arcs $AB$, $BC$ and $CA$ are, respectively, $x+75^{\circ} , 2x+25^{\circ},3x-22^{\circ}$. Then one interior angle of the triangle is:


 $\text{(A) } 57\frac{1}{2}^{\circ}\quad \text{(B) } 59^{\circ}\quad \text{(C) } 60^{\circ}\quad \text{(D) } 61^{\circ}\quad \text{(E) } 122^{\circ}$

 \textbf{Answer: }D



\textbf{Solution}


                    \begin{figure}[H]
                        \centering
                        \includegraphics[width=0.5\textwidth]{Solution/AHSME_I_1969_8_Solution_Figure_1.png}
                    \end{figure}
                    Because the triangle is inscribed, the sum of the minor arcs equals $360^\circ$.  Thus,
 \[x+75+2x+25+3x-22=360\]
 \[6x+78=360\]Solving this yields $x = 47$, so the inscribed angles are $122^\circ$, $99^\circ$, and $119^\circ$.  Noting that an angle of $\triangle ABC$ is half of its corresponding inscribed angle, so the angles of $\triangle ABC$ are $59.5^\circ$, $49.5^\circ$, and $\boxed{\textbf{(D) } 61^\circ}$.

 


	\item (\textit{AHSME 1969}) The arithmetic mean (ordinary average) of the fifty-two successive positive integers beginning at 2 is:


 $\text{(A) } 27\quad \text{(B) } 27\frac{1}{4}\quad \text{(C) } 27\frac{1}{2}\quad \text{(D) } 28\quad \text{(E) } 28\frac{1}{2}$

 \textbf{Answer: }C



\textbf{Solution}

To solve the problem, find the sum of the first $52$ terms of an \href{https://artofproblemsolving.com/wiki/index.php?title=Arithmetic\_sequence}{arithmetic sequence} with first term $2$ and common difference $1$ and divide that by $52$.  The $52^\text{nd}$ term of the sequence is $2+51=53$, so the sum of the first $52$ terms of the sequence is $\frac{52(2+53)}{2} = 1430$.  Thus, the \href{https://artofproblemsolving.com/wiki/index.php?title=Arithmetic\_mean}{arithmetic mean} is $\frac{1430}{52} = \frac{55}{2} = \boxed{\textbf{(C) } 27\frac{1}{2}}$.

 


	\item (\textit{AHSME 1969}) The number of points equidistant from a circle and two parallel tangents to the circle is:


 $\text{(A) } 0\quad \text{(B) } 2\quad \text{(C) } 3\quad \text{(D) } 4\quad \text{(E) } \infty$

 \textbf{Answer: }C



\textbf{Solution}


                    \begin{figure}[H]
                        \centering
                        \includegraphics[width=0.5\textwidth]{Solution/AHSME_I_1969_10_Solution_Figure_1.png}
                    \end{figure}
                    The distance between the two parallel tangents is the length of the circle's diameter, so the distance from a point that satisfies the conditions and the two tangents is the length of the circle's radius.  From the diagram, there are $\boxed{\textbf{(C) } 3}$ points that satisfies the conditions.

 


	\item (\textit{AHSME 1969}) Given points $P(-1,-2)$ and $Q(4,2)$ in the $xy$-plane; point $R(1,m)$ is taken so that $PR+RQ$ is a minimum. Then $m$ equals:


 $\text{(A) } -\frac{3}{5}\quad \text{(B) } -\frac{2}{5}\quad \text{(C) } -\frac{1}{5}\quad \text{(D) } \frac{1}{5}\quad \text{(E) either }-\frac{1}{5}\text{ or} \frac{1}{5}.$

 \textbf{Answer: }B



\textbf{Solution}


                    \begin{figure}[H]
                        \centering
                        \includegraphics[width=0.5\textwidth]{Solution/AHSME_I_1969_11_Solution_Figure_1.png}
                    \end{figure}
                    By the \href{https://artofproblemsolving.com/wiki/index.php?title=Triangle\_Inequality}{Triangle Inequality}, $PR + QR \ge PR$, and equality holds if $R$ is on $PQ$.  The equation of the line with $P$ and $Q$ is $y = \frac{4}{5}x - \frac{6}{5}$, so point $R$ is $(1,-\frac{2}{5})$.  Thus, $m = \boxed{\textbf{(B) } -\frac{2}{5}}$.

 


	\item (\textit{AHSME 1969}) Let $F=\frac{6x^2+16x+3m}{6}$ be the square of an expression which is linear in $x$. Then $m$ has a particular value between:


 $\text{(A) } 3 \text{ and } 4\quad \text{(B) } 4 \text{ and } 5\quad \text{(C) } 5 \text{ and } 6\quad \text{(D) } -4 \text{ and } -3\quad \text{(E) } -6 \text{ and } -5$

 \textbf{Answer: }A



\textbf{Solution}

The quadratic can also be written as
 \[x^2 + \frac{8}{3}x + \frac{m}{2}\]In order for the quadratic to be a square of a linear expression, the constant term must be the square of the x-term divided by 4.  Thus,
 \[\frac{m}{2} = (\frac{8}{3})^2 \cdot \frac{1}{4}\]
 \[\frac{m}{2} = \frac{16}{9}\]
 \[m = \frac{32}{9}\]The answer is $\boxed{\textbf{(A)}}$.

 


	\item (\textit{AHSME 1969}) A circle with radius $r$ is contained within the region bounded by a circle with radius $R$. The area bounded by the larger circle is $\frac{a}{b}$ times the area of the region outside the smaller circle and inside the larger circle. Then $R:r$ equals:


 $\text{(A) }\sqrt{a}:\sqrt{b} \quad \text{(B) } \sqrt{a}:\sqrt{a-b}\quad \text{(C) } \sqrt{b}:\sqrt{a-b}\quad \text{(D) } a:\sqrt{a-b}\quad \text{(E) } b:\sqrt{a-b}$

 \textbf{Answer: }B



\textbf{Solution}

The area of the larger circle is $\pi R^2$, and the area of the region outside the smaller circle and inside the larger circle $\pi R^2 - \pi r^2$.  Thus,
 \[\pi R^2 = \frac{a}{b} \cdot (\pi R^2 - \pi r^2)\]
 \[R^2 = \frac{a}{b} \cdot R^2 - \frac{a}{b} \cdot r^2\]
 \[\frac{a}{b} \cdot r^2 = (\frac{a}{b} -1)R^2\]
 \[\frac{a}{b} \div \frac{a-b}{b} = \frac{R^2}{r^2}\]
 \[\frac{a}{a-b} = \frac{R^2}{r^2}\]
 \[\frac{\sqrt{a}}{\sqrt{a-b}} = \frac{R}{r}\]The answer is $\boxed{\textbf{(B)}}$.

 


	\item (\textit{AHSME 1969}) The complete set of $x$-values satisfying the inequality $\frac{x^2-4}{x^2-1}>0$ is the set of all $x$ such that:


 $\text{(A) } x>2 \text{ or } x<-2 \text{ or} -1<x<1\quad  \text{(B) } x>2 \text{ or } x<-2\quad \\ \text{(C) } x>1 \text{ or } x<-2\qquad\qquad\qquad\quad \text{(D) } x>1 \text{ or } x<-1\quad \\ \text{(E) } x \text{ is any real number except 1 or -1}$

 \textbf{Answer: }A



\textbf{Solution}

Factor the \href{https://artofproblemsolving.com/wiki/index.php?title=Difference\_of\_squares}{difference of squares}.
 \[\frac{(x+2)(x-2)}{(x+1)(x-1)}>0\]Note that the graph intersects the x-axis at when $x = \pm2$ or $x \pm 1$, so check the sign of the result to see if it is positive.  After testing, $x<-2$ or $-1<x<1$ or $x>2$, so the answer is $\boxed{\textbf{(A)}}$.

 


	\item (\textit{AHSME 1969}) In a circle with center $O$ and radius $r$, chord $AB$ is drawn with length equal to $r$ (units). From $O$, a perpendicular to $AB$ meets $AB$ at $M$. From $M$ a perpendicular to $OA$ meets $OA$ at $D$. In terms of $r$ the area of triangle $MDA$, in appropriate square units, is:


 $\text{(A) } \frac{3r^2}{16}\quad \text{(B) } \frac{\pi r^2}{16}\quad \text{(C) } \frac{\pi r^2\sqrt{2}}{8}\quad \text{(D) } \frac{r^2\sqrt{3}}{32}\quad \text{(E) } \frac{r^2\sqrt{6}}{48}$

 \textbf{Answer: }D



\textbf{Solution}


                    \begin{figure}[H]
                        \centering
                        \includegraphics[width=0.5\textwidth]{Solution/AHSME_I_1969_15_Solution_Figure_1.png}
                    \end{figure}
                    

 Because $AO = OB = AB$, $\triangle AOB$ is an \href{https://artofproblemsolving.com/wiki/index.php?title=Equilateral\_triangle}{equilateral triangle}, and $\angle OAB = 60^\circ$.  Using 30-60-90 triangles, $AM = \frac{r}{2}$, $AD = \frac{r}{4}$, and $DM = \frac{r\sqrt{3}}{4}$.  Thus, the area of $\triangle ADM$ is $\frac{1}{2} \cdot \frac{r}{4} \cdot \frac{r\sqrt{3}}{4} = \boxed{\textbf{(D) } \frac{r^2 \sqrt{3}}{32}}$.

 


	\item (\textit{AHSME 1969}) When $(a-b)^n,n\ge2,ab\ne0$, is expanded by the binomial theorem, it is found that when $a=kb$, where $k$ is a positive integer, the sum of the second and third terms is zero. Then $n$ equals:


 $\text{(A) } \frac{1}{2}k(k-1)\quad \text{(B) } \frac{1}{2}k(k+1)\quad \text{(C) } 2k-1\quad \text{(D) } 2k\quad \text{(E) } 2k+1$

 \textbf{Answer: }E



\textbf{Solution}

Since $a=kb$, we can write $(a-b)^n$ as $(kb-b)^n$.Expanding, the second term is $-k^{n-1}b^{n}{{n}\choose{1}}$, and the third term is $k^{n-2}b^{n}{{n}\choose{2}}$, so we can write the equation
 \[-k^{n-1}b^{n}{{n}\choose{1}}+k^{n-2}b^{n}{{n}\choose{2}}=0\]Simplifying and multiplying by two to remove the denominator, we get
 \[-2k^{n-1}b^{n}n+k^{n-2}b^{n}n(n-1)=0\]Factoring, we get
 \[k^{n-2}b^{n}n(-2k+n-1)=0\]Dividing by $k^{n-2}b^{n}$ gives
 \[n(-2k+n-1)=0\]Since it is given that $n\ge2$, $n$ cannot equal 0, so we can divide by n, which gives
 \[-2k+n-1=0\]Solving for $n$ gives
 \[n=2k+1\] so the answer is $\fbox{E}$.

 


	\item (\textit{AHSME 1969}) The equation $2^{2x}-8\cdot 2^x+12=0$ is satisfied by:


 $\text{(A) } log(3)\quad \text{(B) } \frac{1}{2}log(6)\quad \text{(C) } 1+log(\frac{3}{2})\quad \text{(D) } 1+\frac{log(3)}{log(2)}\quad \text{(E) none of these}$

 \textbf{Answer: }D



\textbf{Solution}

Let $2^x=a$. Because $2^{2x}=(2^x)^2$, the given expression can be rewritten as $a^2-8a+12=0$. This can be factored as $(a-6)(a-2)=0$, which has solutions $a=2^x=6$ and $a=2^x=2$. Looking at the answer choices, we see that $x=1$ is absent. Rewriting $2^x=6$ as $x=log_26$ and then applying the logarithm addition identity in reverse gives $x=log_2(2)+log_2(3)=1+log_2(3)$. Applying the logarithm division identity shows that the answer is $\fbox{D}$.

 


	\item (\textit{AHSME 1969}) The number of points common to the graphs of
$(x-y+2)(3x+y-4)=0 \text{    and     } (x+y-2)(2x-5y+7)=0$
is:


 $\text{(A) } 2\quad \text{(B) } 4\quad \text{(C) } 6\quad \text{(D) } 16\quad \text{(E) } \infty$

 \textbf{Answer: }B



\textbf{Solution}

By the \href{https://artofproblemsolving.com/wiki/index.php?title=Zero\_Product\_Property}{Zero Product Property}, $x-y+2=0$ or $3x+y-4=0$ in the first equation and $x+y-2=0$ or $2x-5y+7=0$ in the second equation.  Thus, from the first equation, $y = x+2$ or $y =-3x+4$, and from the second equation, $y=-x+2$ or $y = \frac{2}{5}x + \frac{7}{5}$.

 If a point is common to the two graphs, then the point must be in one of the lines in the first equation as well as one of the lines in the second equation.  Since the slopes of the lines are different, none of the lines are parallel.  Thus, there are $2 \cdot 2 = \boxed{\textbf{(B) } 4}$ points of intersection in the two graphs.

 


	\item (\textit{AHSME 1969}) The number of distinct ordered pairs $(x,y)$ where $x$ and $y$ have positive integral values satisfying the equation $x^4y^4-10x^2y^2+9=0$ is:


 $\text{(A) } 0\quad \text{(B) } 3\quad \text{(C) } 4\quad \text{(D) } 12\quad \text{(E) } \infty$

 \textbf{Answer: }B



\textbf{Solution}

Let $(xy)^2=a$. The expression given is equal to $a^2-10a+9=0$, which can be factored as $(a-9)(a-1)=0$. Thus, we have $a=(xy)^2=9$ and $a=(xy)^2=1$. Because $x$ and $y$ are positive, we can eliminate the possibilities where $xy$ is negative. From here it is easy to see that the only integral pairs of $x$ and $y$ are $(3, 1), (1, 3)$, and $(1, 1)$. The answer is $\fbox{B}$.

 


	\item (\textit{AHSME 1969}) Let $P$ equal the product of 3,659,893,456,789,325,678 and 342,973,489,379,256. The number of digits in $P$ is:


 $\text{(A) } 36\quad \text{(B) } 35\quad \text{(C) } 34\quad \text{(D) } 33\quad \text{(E) } 32$

 \textbf{Answer: }C



\textbf{Solution 1}

Through inspection, we see that the two digit number $33^{2}=1089=4$ digits.Notice that any number that has the form $33abcdefg.......$ multiplied by another $33qwertyu.........$ will have its number of digits equal to the sum of the original numbers' digits.

 In this case, we see that the first number has $19$ digits, and the second number has $15$ digits.

 Note: this applies for numbers $33--->99$

 Hence, the answer is $19+15=34$ digits $\implies \fbox{C}$




\textbf{Solution 2}

We can approximate the product with $10^{18} * 3.6 *10^{14} *3.4$ Now observe that $3.6*3.4>10$, so we can further simplify the product with $10^{18}*10^{14}*10=10^{33}$ which means the product has $\fbox{34 (C)}$ digits.

 -serpent\_121

 


	\item (\textit{AHSME 1969}) If the graph of $x^2+y^2=m$ is tangent to that of $x+y=\sqrt{2m}$, then:


 $\text{(A) m must equal } \frac{1}{2}\quad \text{(B) m must equal  } \frac{1}{\sqrt{2}}\quad\\ \text{(C) m must equal } \sqrt{2}\quad \text{(D) m must equal } 2\quad\\ \text{(E) m may be any non-negative real number}$

 \textbf{Answer: }E



\textbf{Solution 1}

Note that the first equation represents a circle and the second equation represents a line.  If a line is tangent to a circle, then it only hits at one point, so there will only be one solution to the \href{https://artofproblemsolving.com/wiki/index.php?title=System\_of\_equations}{system of equations}.

 In the second equation, $y = -x + \sqrt{2m}$.  Substitution results in 
 \[x^2 + x^2 - 2x\sqrt{2m} + 2m = m\]
 \[2x^2 - 2x\sqrt{2m} + m = 0\]In order for the system to have one solution, the \href{https://artofproblemsolving.com/wiki/index.php?title=Discriminant}{discriminant} must equal $0$.
 \[4(2m) - 4 \cdot 2 \cdot m = 0\]
 \[0 = 0\]Thus, $m$ can be a non-negative real number, so the answer is $\boxed{\textbf{(E)}}$.

 



\textbf{Solution 2}

Since $x^2+y^2=m$ has \href{https://artofproblemsolving.com/wiki/index.php?title=Radius}{radius} $\sqrt{m}$, and $x+y=\sqrt{2m}$ is a \href{https://artofproblemsolving.com/wiki/index.php?title=Diagonal}{diagonal} \href{https://artofproblemsolving.com/wiki/index.php?title=Line}{line} in the \href{https://artofproblemsolving.com/wiki/index.php?title=Plane}{plane} with \href{https://artofproblemsolving.com/wiki/index.php?title=Slope}{slope} $-1$, if the two \href{https://artofproblemsolving.com/wiki/index.php?title=Curves}{curves} are \href{https://artofproblemsolving.com/wiki/index.php?title=Tangent}{tangent} then $x+y=\sqrt{2m}$ either has to pass through $(\frac{\sqrt{2m}}{2}, \frac{\sqrt{2m}}{2})$ or $(-\frac{\sqrt{2m}}{2}, -\frac{\sqrt{2m}}{2})$. However, as $\sqrt{2m}$ is \href{https://artofproblemsolving.com/wiki/index.php?title=Positive}{positive}, the \href{https://artofproblemsolving.com/wiki/index.php?title=Line}{line} can only pass through $(\frac{\sqrt{2m}}{2}, \frac{\sqrt{2m}}{2})$, which it does for all $m$. Thus the answer is $\boxed{\textbf{(E)}}$.

 


	\item (\textit{AHSME 1969}) Let $K$ be the measure of the area bounded by the $x$-axis, the line $x=8$, and the curve defined by


 
 \[f={(x,y)\quad |\quad y=x \text{  when  } 0 \le x \le 5, y=2x-5 \text{  when  } 5 \le x \le 8}.\]


 Then $K$ is:


 $\text{(A) } 21.5\quad \text{(B) } 36.4\quad \text{(C) } 36.5\quad \text{(D) } 44\quad \text{(E) less than 44 but arbitrarily close to it}$

 \textbf{Answer: }C



\textbf{Solution}


                    \begin{figure}[H]
                        \centering
                        \includegraphics[width=0.5\textwidth]{Solution/AHSME_I_1969_22_Solution_Figure_1.png}
                    \end{figure}
                    

 The shape can be divided into a triangle and a trapezoid.  For the triangle, the base is $5$ and the height is $5$, so the area is $\frac{5 \cdot 5}{2} = \frac{25}{2}$.  For the trapezoid, the two bases are $5$ and $11$ and the height is $3$, so the area is $\frac{3(5+11)}{2} = 24$.  Thus, the total area is $\frac{25}{2} + 24 = \frac{73}{2} = \boxed{\textbf{(C) } 36.5}$.

 


	\item (\textit{AHSME 1969}) For any integer $n>1$, the number of prime numbers greater than $n!+1$ and less than $n!+n$ is:


 $\text{(A) } 0\quad\qquad \text{(B) } 1\quad\\ \text{(C) } \frac{n}{2} \text{ for n even, } \frac{n+1}{2} \text{ for n odd}\quad\\ \text{(D) } n-1\quad \text{(E) } n$

 \textbf{Answer: }A



\textbf{Solution}

Observe that for all $k \in 1< k< n$, since $k$ divides $n!$, $k$ also divides $n!+k$. Therefore, all numbers $a$ in the range $n!+1<a<n!+n$ are composite. Therefore there are 0 primes in that range. $\fbox{A}$



	\item (\textit{AHSME 1969}) When the natural numbers $P$ and $P'$, with $P>P'$, are divided by the natural number $D$, the remainders are $R$ and $R'$, respectively. When $PP'$ and $RR'$ are divided by $D$, the remainders are $r$ and $r'$, respectively. Then:


 $\text{(A) } r>r' \text{  always}\quad \text{(B) } r<r' \text{  always}\quad\\ \text{(C) } r>r' \text{  sometimes and } r<r' \text{  sometimes}\quad\\ \text{(D) } r>r' \text{  sometimes and } r=r' \text{  sometimes}\quad\\ \text{(E) } r=r' \text{  always}$

 \textbf{Answer: }E



\textbf{Solution}

The divisors are the same, so take each variable \href{https://artofproblemsolving.com/wiki/index.php?title=Modulo}{modulo} $D$.
 \[P \equiv R \pmod{D}\]
 \[P' \equiv R’ \pmod{D}\]That means
 \[PP’ \equiv RR’ \pmod{D}\]Thus, $PP’$ and $RR’$ have the same remainder when divided by $D$, so the answer is $\boxed{\textbf{(E)}}$.

 


	\item (\textit{AHSME 1969}) If it is known that $\log_2(a)+\log_2(b) \ge 6$, then the least value that can be taken on by $a+b$ is:


 $\text{(A) } 2\sqrt{6}\quad \text{(B) } 6\quad \text{(C) } 8\sqrt{2}\quad \text{(D) } 16\quad \text{(E) none of these}$

 \textbf{Answer: }D



\textbf{Solution}

We use the logarithm property of addition:
 \begin{align*} \log_2(a)+\log_2(b) \ge 6 &= \log_2(ab) \ge 6\\ &\Rightarrow 2^{log_2(ab)} \ge 2^6\\ &= ab \ge 64 \end{align*}
Due to the Quadratic \href{https://artofproblemsolving.com/wiki/index.php?title=Optimization}{Optimization} or the \href{https://artofproblemsolving.com/wiki/index.php?title=AM-GM\_Inequality}{AM-GM Inequality}, the least value is obtained when $a = b$.Therefore, $a = b = 8 \Rightarrow a + b = \boxed{(D)16}$



	\item (\textit{AHSME 1969}) 
                    \begin{figure}[H]
                        \centering
                        \includegraphics[width=0.5\textwidth]{Problem/AHSME_I_1969_26_Problem_Figure_1.png}
                    \end{figure}
                    


 A parabolic arch has a height of $16$ inches and a span of $40$ inches. The height, in inches, of the arch at the point $5$ inches from the center $M$ is:


 $\text{(A) } 1\quad \text{(B) } 15\quad \text{(C) } 15\frac{1}{3}\quad \text{(D) } 15\frac{1}{2}\quad \text{(E) } 15\frac{3}{4}$

 \textbf{Answer: }B



\textbf{Solution}

Because the arch has a height of $16$ inches, an equation that models the arch is $y = ax^2 + 16$, where $x$ is the horizontal distance from the center and $y$ is the height.  The arch has a span of $40$ inches, so the arch meets the ground $20$ inches from the center.  That means $0 = 400a + 16$, so $a = -\frac{1}{25}$.

 Thus, the equation that models height based on distance from the center $y = -\frac{1}{25}x^2 + 16$, so the height of the arch $5$ inches from the center is $-\frac{1}{25} \cdot 5^2 + 16 = \boxed{\textbf{(B) } 15}$ inches.

 


	\item (\textit{AHSME 1969}) A particle moves so that its speed for the second and subsequent miles varies inversely as the integral number of miles already traveled. For each subsequent mile the speed is constant. If the second mile is traversed in $2$ hours, then the time, in hours, needed to traverse the $n$th mile is:


 $\text{(A) } \frac{2}{n-1}\quad \text{(B) } \frac{n-1}{2}\quad \text{(C) } \frac{2}{n}\quad \text{(D) } 2n\quad \text{(E) } 2(n-1)$

 \textbf{Answer: }E



\textbf{Solution}

Let $d$ be the distance already traveled and $s_n$ be the speed for the $n^\text{th}$ mile.  Because the speed for the second and subsequent miles varies inversely as the integral number of miles already traveled,
 \[ds_{d+1} = k\]If the second mile is traversed in $2$ hours, then the speed in the second mile ($s_2$) equals $\frac{1}{2}$ miles per hour.  Since one mile has been traveled already, $k = \frac{1}{2}$.  Now solve for the speed the $n^\text{th}$ mile, noting that $n-1$ miles have already been traveled.
 \[(n-1)s_{d+1} = \frac{1}{2}\]
 \[s_{d+1} = \frac{1}{2(n-1)}\]Thus, the time it takes to travel the $n^\text{th}$ mile is $1 \div \frac{1}{2(n-1)} = \boxed{\textbf{(E) } 2(n-1)}$ hours.

 


	\item (\textit{AHSME 1969}) Let $n$ be the number of points $P$ interior to the region bounded by a circle with radius $1$, such that the sum of squares of the distances from $P$ to the endpoints of a given diameter is $3$. Then $n$ is:


 $\text{(A) } 0\quad \text{(B) } 1\quad \text{(C) } 2\quad \text{(D) } 4\quad \text{(E) } \infty$

 \textbf{Answer: }E



\textbf{Solution}


                    \begin{figure}[H]
                        \centering
                        \includegraphics[width=0.5\textwidth]{Solution/AHSME_I_1969_28_Solution_Figure_1.png}
                    \end{figure}
                    

 Let $A$ and $B$ be points on diameter.  Extend $AP$, and mark intersection with circle as point $C$.

 Because $AB$ is a diameter, $\angle ACB = 90^\circ$.  Also, by Exterior Angle Theorem, $\angle ACB + \angle CBP = \angle APB$, so $\angle APB > \angle ACB$, making $\angle APB$ an obtuse angle.

 By the \href{https://artofproblemsolving.com/wiki/index.php?title=Law\_of\_Cosines}{Law of Cosines}, $AP^2 + BP^2 - 2 \cdot AP \cdot BP \cdot \cos{\angle APB} = 4$.  Since $AP^2 + BP^2 = 3$, substitute and simplify to get $\cos{\angle APB} = -\frac{1}{2 \cdot AP \cdot BP}$.  This equation has infinite solutions because for every $AP$ and $BP$, where $AP + BP \ge 2$ and $AP$ and $BP$ are both less than $2$, there can be an obtuse angle that satisfies the equation, so the answer is $\boxed{\textbf{(E)}}$.

 



\textbf{Solution 2 (Apollonius Theorem)}

Let $C_1$ be the circle with a centra $O$ and diameter $AB$. Let $P$ be any random point inside $C_1$. Connect $OP$. We want $AP^2 + BP^2 = 3$In $\triangle APB, PO$ is a median on $AB$. Using Appolonius Theorem we get:
 \[AP^2 + BP^2 = 2(AO^2 + OP^2)\]
 \[2(OP^2 + 1) = 3\]
 \[2\cdot OP^2 + 2 = 3\]
 \[OP = \frac{1}{\sqrt{2}}\]

 We draw a circle $C_2$ centred at $O$ with a radius of $\frac{1}{\sqrt{2}}$ i.e, $OP$.Every point on the circle $C_2$ will satisfy the given conditions. Since there are infinite points on a circle hence there are infinite number of points P that can satisfy the condition. 

 Hence, the answer is $\boxed{\textbf{(E)}}$.

 --Tanstorm

 


	\item (\textit{AHSME 1969}) If $x=t^{1/(t-1)}$ and  $y=t^{t/(t-1)},t>0,t \ne 1$, a relation between $x$ and $y$ is: 


 $\text{(A) } y^x=x^{1/y}\quad \text{(B) } y^{1/x}=x^{y}\quad \text{(C) } y^x=x^y\quad \text{(D) } x^x=y^y\quad \text{(E) none of these}$

 \textbf{Answer: }C



\textbf{Solution}

Plug in $t = 2$ and test each expression (if the relation works, then both sides are equal).  After testing, options A, B, and D are out, but for option C, both sides are equal.  To check that C is a valid option, substitute the values of $x$ and $y$ and use exponent properties.
 \[x^y = (t^{\frac{1}{t-1}})^{t^{\frac{t}{t-1}}} = t^{\frac{1}{t-1} \cdot t^{\frac{t}{t-1}}}\]
 \[y^x = (t^{\frac{t}{t-1}})^{t^{\frac{1}{t-1}}} = t^{\frac{t}{t-1} \cdot t^{\frac{1}{t-1}}}\]In the second equation, the exponent can be rewritten as $\frac{1}{t-1} \cdot t \cdot t^{\frac{1}{t-1}}$.  Since $\frac{t-1}{t-1} = 1$, the exponent can be simplified as $\frac{1}{t-1} \cdot t^{\frac{t-1}{t-1}} \cdot t^{\frac{1}{t-1}} = \frac{1}{t-1} \cdot t^{\frac{t}{t-1}}$.  Since the base and exponent of both $x^y$ and $y^x$ are the same, we can confirm that $\boxed{\textbf{(C) } y^x = x^y}$.

 


	\item (\textit{AHSME 1969}) Let $P$ be a point of hypotenuse $AB$ (or its extension) of isosceles right triangle $ABC$. Let $s=AP^2+PB^2$. Then:


 $\text{(A) } s<2CP^2 \text{ for a finite number of positions of P}\quad\\ \text{(B) } s<2CP^2 \text{ for an infinite number of positions of P}\quad\\ \text{(C) } s=2CP^2 \text{ only if P is the midpoint or an endpoint of AB}\quad\\ \text{(D) } s=2CP^2 \text{ always}\quad\\ \text{(E) } s>2CP^2 \text{ if P is a trisection point of AB}$

 \textbf{Answer: }D



\textbf{Solution}


                    \begin{figure}[H]
                        \centering
                        \includegraphics[width=0.5\textwidth]{Solution/AHSME_I_1969_30_Solution_Figure_1.png}
                    \end{figure}
                    

 Consider the case where $P$ is on the hypotenuse of $AB$.  Draw perpendicular lines from $P$ towards the sides.  Using the \href{https://artofproblemsolving.com/wiki/index.php?title=Pythagorean\_Theorem}{Pythagorean Theorem}, 
 \[AP^2 = a^2 + a^2\]
 \[BP^2 = (x-a)^2 + (x-a)^2\]
 \[CP^2 = a^2 + (x-a)^2\]This means
 \[s = 4a^2 - 4ax + 2x^2\]
 \[2 \cdot CP^2 = 4a^2 - 4ax + 2x^2\]Thus, $s = 2 \cdot CP^2$ when $P$ is on the hypotenuse of $AB$.

 

                    \begin{figure}[H]
                        \centering
                        \includegraphics[width=0.5\textwidth]{Solution/AHSME_I_1969_30_Solution_Figure_2.png}
                    \end{figure}
                    

 Consider the case where $P$ is on the extension of $AB$.  \href{https://artofproblemsolving.com/wiki/index.php?title=WLOG}{WLOG}, let point $A$ be between point $P$ and point $B$.  Extend $BC$ and draw perpendicular line from $P$.  Also, draw point $D$, where $PD \parallel AC$ and $DA \parallel CB$.

 Using the Pythagorean Theorem again,
 \[AP^2 = a^2 + a^2\]
 \[BP^2 = (a+x)^2 + (a+x)^2\]
 \[CP^2 = (a+x)^2 + a^2\]That means
 \[s = 4a^2 + 4ax + 2x^2\]
 \[2 \cdot CP^2 = 4a^2 + 4ax + 2x^2\]Thus, $s = 2 \cdot CP^2$ when $P$ is outside the hypotenuse.

 In summary, $AP^2 + BP^2 = 2 \cdot CP^2$, so the answer is $\boxed{\textbf{(D)}}$.

 


	\item (\textit{AHSME 1969}) Let $OABC$ be a unit square in the $xy$-plane with $O(0,0),A(1,0),B(1,1)$ and $C(0,1)$. Let $u=x^2-y^2$, and $v=2xy$ be a transformation of the $xy$-plane into the $uv$-plane. The transform (or image) of the square is:


 
                    \begin{figure}[H]
                        \centering
                        \includegraphics[width=0.5\textwidth]{Problem/AHSME_I_1969_31_Problem_Figure_1.png}
                    \end{figure}
                    

                    \begin{figure}[H]
                        \centering
                        \includegraphics[width=0.5\textwidth]{Problem/AHSME_I_1969_31_Problem_Figure_2.png}
                    \end{figure}
                    

                    \begin{figure}[H]
                        \centering
                        \includegraphics[width=0.5\textwidth]{Problem/AHSME_I_1969_31_Problem_Figure_3.png}
                    \end{figure}
                    

                    \begin{figure}[H]
                        \centering
                        \includegraphics[width=0.5\textwidth]{Problem/AHSME_I_1969_31_Problem_Figure_4.png}
                    \end{figure}
                    

                    \begin{figure}[H]
                        \centering
                        \includegraphics[width=0.5\textwidth]{Problem/AHSME_I_1969_31_Problem_Figure_5.png}
                    \end{figure}
                    

 \textbf{Answer: }D



\textbf{Solution}

Each point on the square can be in the form $(0,y)$, $(1,y)$, $(x,0)$, and $(x,1)$, where $0 \le x,y \le 1$.  Making the appropriate substitutions results in points being $(-y^2,0)$, $(1-y^2,2y)$, $(x^2,0)$, and $(x^2 - 1,2x)$ on the $uv$-plane.  

 Notice that since $v \ge 0$, none of the points are below the u-axis, so options A,B, and E are out.  Since $x = \frac{v}{2}$, $u = (\frac{v}{2})^2 - 1$, so $v = 2\sqrt{u+1}$, where $-1 \le u \le 0$.  That means some of the lines are not straight, so the answer is $\boxed{\textbf{(D)}}$.

 


	\item (\textit{AHSME 1969}) Let a sequence $\{u_n\}$ be defined by $u_1=5$ and the relationship $u_{n+1}-u_n=3+4(n-1), n=1,2,3\cdots.$If $u_n$ is expressed as a polynomial in $n$, the algebraic sum of its coefficients is:


 $\text{(A) 3} \quad \text{(B) 4} \quad \text{(C) 5} \quad \text{(D) 6} \quad \text{(E) 11}$

 \textbf{Answer: }C



\textbf{Solution}

Note that the first differences create a linear function, so the sequence ${u_n}$ is quadratic.

 The first three terms of the sequence are $5$, $8$, and $15$.  From there, a \href{https://artofproblemsolving.com/wiki/index.php?title=System\_of\_equations}{system of equations} can be written.
 \[a+b+c=5\]
 \[4a+2b+c=8\]
 \[9a+3b+c=15\]Solve the system to get $a=2$, $b=-3$, and $c=6$.  The sum of the coefficients is $\boxed{\textbf{(C) } 5}$.

 \begin{verbatim}
Note: Solving the system is extra work, as the answer is described by the first equation. The sum of the coefficients ($a + b + c$) is just 5 by the first equation.
\end{verbatim}




	\item (\textit{AHSME 1969}) Let $S_n$ and $T_n$ be the respective sums of the first $n$ terms of two arithmetic series. If $S_n:T_n=(7n+1):(4n+27)$ for all $n$, the ratio of the eleventh term of the first series to the eleventh term of the second series is:


 $\text{(A) } 4:3\quad \text{(B) } 3:2\quad \text{(C) } 7:4\quad \text{(D) } 78:71\quad \text{(E) undetermined}$

 \textbf{Answer: }A



\textbf{Solution}

Let $S$ be the first arithmetic sequence and $T$ be the second arithmetic sequence.  If $n = 1$, then $S_1:T_1 = 8:31$.  Since $S_1$ and $T_1$ are just the first term, the first term of $S$ is $8a$ and the first term of $T$ is $31a$ for some $a$.  If $n = 2$, then $S_2:T_2 = 15:35 = 3:7$, so the sum of the first two terms of $S$ is $3b$ and the sum of the first two terms of $T$ is $7b$ for some $b$.  Thus, the second term of $S$ is $3b-8a$ and the second term of $T$ is $7b - 31a$, so the common difference of $S$ is $3b-16a$ and the common difference of $T$ is $7b-62a$.

 Thus, using the first terms and common differences, the sum of the first three terms of $S$ equals $\frac{1}{2} \cdot 3(16a + 2(-16a + 3b))$, and the sum of the first three terms of $T$ equals $\frac{1}{2} \cdot 3(62a + 2(-62a + 7b))$.  That means
 \[\frac{16a + 2(-16a + 3b)}{62a + 2(-62a + 7b)} = \frac{22}{39}\]
 \[\frac{6b-16a}{14b-62a} = \frac{22}{39}\]
 \[\frac{3b-8a}{7b-31a} = \frac{22}{39}\]
 \[117b - 312a = 154b - 682a\]
 \[-37b = -370a\]
 \[b = 10a\]With the substitution, the common difference of $S$ is $14a$, and the common difference of $T$ is $8a$.  That means the $11^\text{th}$ term of $S$ is $8a + 10(14a) = 148a$, and the $11^\text{th}$ term of $T$ is $31a + 10(8a) = 111a$.  Thus, the ratio of the eleventh term of the first series to the eleventh term of the second series is $148a:111a = \boxed{\textbf{(A) } 4:3}$.

 



\textbf{Solution 2 (Quick Sol)}

Note that the sum of arithmetic sequences can be expressed as a quadratic polynomial of $n$. From the ratio given in the problem, multiply $n$ to both sides to get quadratic polynomials. 
 \[S_n = 7n^2+n, T_n = 4n^2+27n\] From there, the 11th term for $S_n$ and $T_n$ can be calculated. 
 \[S_{11} - S_{10} = 7*11^2+11 - (7*10^2+10) = 148\] 
 \[T_{11} - T_{10} = 4*11^2+27*11 - (4*10^2+27*10) = 111\] The ratio is $148 : 111 = \boxed{\textbf{(A) } 4:3}$.

 



\textbf{Solution 3 (Quickest Sol)}

Since they are arithmetic sequences, we know that the 11th term is the average value of the first 21 terms in each sequence. So the desired ratio of the 11th terms is just $(21 \cdot 7 + 1) : (21 \cdot 4 + 27) = 148 : 111 = \boxed{\textbf{(A) } 4:3}$

 -purplepenguin2

 



\textbf{Solution 4 (Formal Quick)}

Let a, and c be the first term of $S_n$ and $T_n$ respectively and d, b be the common difference of $S_n$ and $T_n$ respectively


 $\frac{S_n}{T_n} = \frac{\frac{n}{2}(2a + (n-1)d)}{\frac{n}{2}(2c + (n-1)b)} = \frac{7n+1}{4n+27}$
$\frac{S_n}{T_n} = \frac{(2a + (n-1)d)}{(2c + (n-1)b)}$


 Now major step!!!!


 Divide numerator and denominator of R.H.S by 2


 $\frac{S_n}{T_n} = \frac{(a + (\frac{n-1}{2})d)}{(c + (\frac{n-1}{2})b)} = \frac{7n+1}{4n+27}$


 Since we want ratio of 11th terms 


 $\frac{n-1}{2} = 10$
$n = 21$
$\frac{S_{21}}{T_{21}} = \frac{(a + (10)d)}{(c + (10)b)} = \frac{148}{111}$


 which is equal to $\frac{4}{3}$


 Actually this solution was given by my maths teacher.

 


	\item (\textit{AHSME 1969}) The remainder $R$ obtained by dividing $x^{100}$ by $x^2-3x+2$ is a polynomial of degree less than $2$. Then $R$ may be written as:


 

$\text{(A) }2^{100}-1 \quad \text{(B) } 2^{100}(x-1)-(x-2)\quad \text{(C) } 2^{200}(x-3)\quad\\ \text{(D) } x(2^{100}-1)+2(2^{99}-1)\quad \text{(E) } 2^{100}(x+1)-(x+2)$

 \textbf{Answer: }B



\textbf{Solution}

Let the polynomial $Q(x)$ be the quotient when $x^{100}$ is divided by $x^2-3x+2$, and let the remainder $R=ax+b$, for some real $a$ and $b$. Then we can write: $x^{100}=(x^2-3x+2)Q(x)+ax+b$. Since it is hard to deal with $Q(x)$ (it is of degree 98!), we factor $x^2-3x+2$ as $(x-2)(x-1)$ so we can eliminate $Q(x)$ by plugging in $x$ values of $2$ and $1$.

 
$x^{100}=(x-2)(x-1)Q(x)+ax+b$,

 $2^{100}=(0)Q(2)+2a+b$,

 $2^{100}=2a+b$.

 Similarly, $1^{100}=a+b$.

 
Solving this system of equations gives $a=2^{100}-1$ and $b=2-2^{100}$. Thus, $R=ax+b=(2^{100}-1)x+(2-2^{100})$. Expanding and combining terms, we see that the answer is $\fbox{B}$.

 


	\item (\textit{AHSME 1969}) Let $L(m)$ be the $x$ coordinate of the left end point of the intersection of the graphs of $y=x^2-6$ and $y=m$, where $-6<m<6$. Let $r=[L(-m)-L(m)]/m$. Then, as $m$ is made arbitrarily close to zero, the value of $r$ is:


 $\text{(A) arbitrarily close to } 0\quad\\ \text{(B) arbitrarily close to }\frac{1}{\sqrt{6}}\quad \text{(C) arbitrarily close to }\frac{2}{\sqrt{6}} \quad \text{(D) arbitrarily large} \quad \text{(E) undetermined}$

 \textbf{Answer: }B



\textbf{Solution 1}

Since $L(a)$ is the $x$ coordinate of the left end point of the intersection of the graphs of $y=x^2-6$ and $y=a$, we can substitute $a$ for $y$ and find the lowest solution $x$.
 \[x^2 - 6 = a\]
 \[x = \pm \sqrt{a+6}\]That means $L(m) = -\sqrt{m+6}$ and $L(-m) = -\sqrt{-m+6}$.  That means
 \[r = \frac{-\sqrt{-m+6}+\sqrt{m+6}}{m}\]Since plugging in $0$ for $m$ results in $0/0$, there is a removable discontinuity.  Multiply the fraction by $\frac{\sqrt{-m+6}+\sqrt{m+6}}{\sqrt{-m+6}+\sqrt{m+6}}$ to get
 \[r = \frac{m+6 - (-m+6)}{m(\sqrt{-m+6}+\sqrt{m+6})}\]
 \[r = \frac{2m}{m(\sqrt{-m+6}+\sqrt{m+6})}\]
 \[r = \frac{2}{(\sqrt{-m+6}+\sqrt{m+6})}\]Now there wouldn't be a problem plugging in $0$ for $m$.  Doing so results in $r = \frac{2}{2\sqrt{6}} = \frac{1}{\sqrt{6}}$, so the answer is $\boxed{\textbf{(B)}}$.

 



\textbf{Solution 2}

From Solution 1, $L(m) = -\sqrt{m+6}$ and $L(-m) = -\sqrt{-m+6}$, so
 \[r = \frac{-\sqrt{-m+6}+\sqrt{m+6}}{m}\]Since $m$ is arbitrarily close to $0$, we wish to find
 \[\lim_{m \rightarrow 0} \frac{-\sqrt{-m + 6} + \sqrt{m + 6}}{m}\]Using \href{https://artofproblemsolving.com/wiki/index.php?title=L\%27Hopital\%27s\_Rule}{L'Hopital's Rule}, the limit is equivalent to
 \[\lim_{m \rightarrow 0} \frac{1}{2 \sqrt{-m + 6}} + \frac{1}{2 \sqrt{m + 6}}\]Calculating the limit shows that $r$ is $\boxed{\text{(B) arbitrarily close to }\frac{1}{\sqrt{6}}}$.

 


\end{enumerate}